\providecommand\ExamFontPath{../../../字体/}
\documentclass[fontset=exam]{exam-zh}
\usepackage{longtable,booktabs,array,calc,multirow,fvextra}
\usepackage[normalem]{ulem}
\providecommand{\hl}[1]{#1}
\providecommand{\ul}[1]{\uline{#1}}
\DefineVerbatimEnvironment{verbatim}{Verbatim}{breaklines=true,fontsize=\small}
\providecommand{\tightlist}{\setlength{\itemsep}{0pt}\setlength{\parskip}{0pt}}
\newcommand{\pandocbounded}[1]{#1}
\setlength{\emergencystretch}{3em}
\setlength{\tabcolsep}{2pt}
\examsetup{page/size=a4paper}
\title{厦门大学 2020 级 期中考试 A 卷}
\subject[4em]{概率统计}
\begin{document}
\maketitle
\fancyfoot[C]{第 \thepage 页（共 \pageref*{LastPage} 页）}
\pdfbookmark[1]{一、选择题（在各小题四个备选答案中选出一个正确答案，填在题后的括号中，本大题共6个小题，每小题3分，总计18分）}{section4}

\par\medskip{\heiti

一、选择题（在各小题四个备选答案中选出一个正确答案，填在题后的括号中，本大题共 6 个小题，每小题 3 分，总计 18 分）

\par}\smallskip

1. 连续型随机变量\(\left. X \right.\sim N(0,1)\)，其分布函数为\(\Phi(x) = \frac{1}{\sqrt{2\pi}}{\int_{- \infty}^{x}e^{- \frac{x^{2}}{2}}}dt\left( {x \geq 0} \right)\)，则下式成立的是（ ）。

(A) \(P\left( \left| x \middle| > 2 \right. \right) = 2\lbrack 1 - \Phi(2)\rbrack\) (B) \(\Phi(x) = 1 + \Phi( - x)\)

(C) \(\Phi(kx) = k\Phi(x)\) (D) \(\Phi(x_{1} + x_{2}) = \Phi(x_{1}) + \Phi(x_{2})\)

2. 设随机变量\(X\)的概率密度函数为\(f(x)\)，则\(Y = 5 - 2X\)的概率密度函数为（ ）。

(A) \(- \frac{1}{2}f( - \frac{y - 5}{2})\) (B) \(\frac{1}{2}f( - \frac{y - 5}{2})\)

(C) \(- \frac{1}{2}f( - \frac{y + 5}{2})\) (D) \(\frac{1}{2}f( - \frac{y + 5}{2})\)

3. 如果\(X,Y\)为两个随机变量，满足\(\rho_{xy} = 0\)，下列命题中错误的是（ ）。

(A) \(X,Y\)不相关 (B) \(X,Y\text{相互独立}\)

(C) \(E\left( {XY} \right) = E(X)E(Y)\) (D) \(D\left( {X - Y} \right) = D(X) + D(Y)\)

4. 设\emph{A}, \emph{B}为随机事件，则 P(A) = P(B)的充分必要条件是（ ）。

\begin{quote}
(A) \(P(A \cup B) = P(A) + P(B)\) (B) \(P(AB) = P(A)P(B)\)

(C) \(P(A\overline{B}) = P(B\overline{A})\) (D) \(P(AB) = P(\overline{A}\overline{B})\)
\end{quote}

5. 设某公路上经过的货车与客车的数量之比为 2：1，货车中途停车修理的概率为 0.02，客车为 0.01，今有一辆汽车中途停车修理，那么该汽车是货车的概率为（ ）。

\begin{quote}
(A) 4/9 (B) 3/5 (C) 2/3 (D) 4/5
\end{quote}

6. 设\(X\text{\textasciitilde}U\left( {a,b} \right)\)，\(E(X) = 3\)，\(D(X) = \frac{1}{3}\)，则\(P\left( {1 < X < 3} \right)\)=（ ）。

(A) 0 (B) \(\frac{1}{2}\) (C) \(\frac{1}{3}\) (D) \(\frac{1}{4}\)

\pdfbookmark[1]{二、填空题（本大题共6小题，每小题3分，第1题第一空1分，第二空2分，总计18分）}{section20}

\par\medskip{\heiti

二、填空题（本大题共 6 小题，每小题 3 分，第 1 题第一空 1 分，第二空 2 分，总计 18 分）

\par}\smallskip

\begin{enumerate}
\def\labelenumi{\arabic{enumi}.}
\item
  将 52 张扑克牌扣在桌子上，一张张地翻开，直到出现一张红桃 K，问接下来的一张牌是黑桃 K 的概率是 ，是红桃 1 的概率是 。
\item
  设随机变量\(X\)服从泊松分布，且\(P(X \leq 1) = 4P(X = 2)\)，则\(P(X = 3) =\) 。
\item
  随机变量\(X\text{\textasciitilde}N(\frac{1}{4},9)\)，以\(Y\)表示对\(X\)的 5 次独立重复观察中``\(X \leq \frac{1}{4}\)''出现的次数，则\(P\left\{ Y = 2 \right\} =\) 。
\item
  设 A, B 独立，A 发生 B 不发生的概率与 A 不发生 B 发生的概率相等，且 A，B 都不发生的概率为\(\frac{1}{9}\)，则\(P(A)\text{=}\) 。
\item
  设事件 A, B 满足条件\(\overline{A} \supset \overline{B}\)，且\(P(A) = \frac{1}{3}\text{，}\allowbreak P(B) = \frac{1}{2}\)，则\(P(\overline{A}B) =\) 。
\item
  设\(E(X) = - 2\text{，}\allowbreak E(Y) = 2\text{，}\allowbreak D(X) = 1\text{，}\allowbreak D(Y) = 4\text{，}\allowbreak \rho_{XY} = - 0.5\)，用切比雪夫不等式估计\(\text{\underline{\hspace{2cm}}}\) \(P\left\{ \middle| X + Y \middle| \geq 6 \right\} \leq\) 。
\end{enumerate}

\pdfbookmark[1]{三、（10分）设有来自三个地区的各10、15、25名考生的报名表，其中女生的报名表分别为3、7、5份，随机取出一个地区的}{section22}

三、（10 分）设有来自三个地区的各 10、15、25 名考生的报名表，其中女生的报名表分别为 3、7、5 份，随机取出一个地区的报名表，从中先后任意抽取两份。求：

（1） 先抽取到的一份是女生表格的概率；

（2） 已知后取到的一份是男生表，先抽取到的一份是女人表的概率。

\pdfbookmark[1]{四、（10分）某射手有3发子弹，射一次命中的概率为frac23，如果命中了就停止设计，否则一直独立射到子弹用尽。求： （}{section25}

四、（10 分）某射手有 3 发子弹，射一次命中的概率为\(\frac{2}{3}\)，如果命中了就停止设计，否则一直独立射到子弹用尽。求：\\
（1） 耗用子弹数\(X\)的分布律； （2） 求\(E(X)\)、\(D(X)\)。

\pdfbookmark[1]{五、（8分）某施工队完成某工程所需要时间X（天）服从Nleft( 100，5^2 right)，施工队上级规定：若工程在}{section26}

五、（8 分）某施工队完成某工程所需要时间\(X\)（天）服从\(N\left( {100\text{，}\allowbreak 5^{2}} \right)\)，施工队上级规定：若工程在 100 天内完成，可以得到奖金 10 万元；在 100\textasciitilde115 天内完成，可以得到奖金 3 万元，若超过 115 天完成，罚款 5 万元。求该施工队在完成该工程时，获得金额的期望。

（附标准正态分布函数值：\(\Phi(1)\text{=}0.8413\text{，}\allowbreak \Phi(2)\text{=}0.97725,\Phi(3) = 0.99865\)）

\pdfbookmark[1]{六、（12分）已知二维随机变量left( X,Y right)的联合概率密度为}{section28}

六、（12 分）已知二维随机变量\(\left( {X,Y} \right)\)的联合概率密度为

\[f_{XY}(x,y) = \left\{ \begin{matrix}
{Ke^{- (2x + y)}} & {x > 0,y > 0} \\
0 & {Others}
\end{matrix} \right.\]

求：

\begin{enumerate}
\def\labelenumi{\arabic{enumi}.}
\item
  边缘概率密度\(f_{X}(x)\)；
\item
  概率\(P\left\{ {Y \leq X} \right\}\)；
\item
  联合分布函数\(F_{XY}(x,y)\)和边缘分布函数\(F_{X}(x)\)。
\end{enumerate}

\pdfbookmark[1]{七、（8分）设随机变量 X 与 Y 相互独立，X的密度函数为f(x) = left beginmatrix
2x，0 <}{section32}

七、（8 分）设随机变量 \emph{X} 与 \emph{Y} 相互独立，\emph{X}的密度函数为\(f(x) = \left\{ \begin{matrix}
{2x\text{，}\allowbreak 0 < x < 1} \\
{0\text{，其他}}
\end{matrix} \right.\)，\(\left. Y \right.\sim B(1,\frac{1}{2})\)，求\(Z = max\left\{ X,Y \right\}\)的分布函数\(F_{Z}(z)\)，并问\emph{Z}是否为连续型随机变量？

\pdfbookmark[1]{八、（8分）随机变量X的概率密度函数为f(x)，且f(x) = f( - x)，F(x)是X的分布函数，证明：对任意实数}{section33}

八、（8 分）随机变量\(X\)的概率密度函数为\(f(x)\)，且\(f(x) = f( - x)\)，\(F(x)\)是\(X\)的分布函数，证明：对任意实数\(a\)，都有 \(F( - a) = \frac{1}{2} - {\int_{0}^{a}{f(x)dx}}\)。

\pdfbookmark[1]{九、（8分）设A，B是两个随机事件，随机变量X = left beginarrayl
1，若A出现
- 1，若A不出现}{section34}

九、（8 分）设\(A\)，\(B\)是两个随机事件，随机变量\(X = \left\{ \begin{array}{l}
{1\text{，若}A\text{出现}} \\
{- 1\text{，若}A\text{不出现}}
\end{array} \right.\)，\(Y = \left\{ \begin{array}{l}
{1,\text{若}B\text{出现}} \\
{\text{-}1,\text{若}B\text{不出现}}
\end{array} \right.\)，证明：随机变量\(X\)和\(Y\)不相关的充要条件是\(A\)和\(B\)相互独立。

\end{document}
